%%%%%%%%%%%%%%%%%%%%%%%%%%%%%%%%%%%%%%%%%%%%%%%%%%%%%%%%%%%%%%%%%%%%%%
% Lớp 12 | Các số đặc trưng ... ghép nhóm | Xác suất có điều kiện
% Dạng 2D6?1-1: Công thức lý thuyết
% Số câu: 20  (sinh tự động bởi tach_cau.py; từ nay thầy sửa trực tiếp file này)
%%%%%%%%%%%%%%%%%%%%%%%%%%%%%%%%%%%%%%%%%%%%%%%%%%%%%%%%%%%%%%%%%%%%%%

%%%%%%%%%%%%%%%% TRẮC NGHIỆM 4 PHƯƠNG ÁN %%%%%%%%%%%%%%%%

% nguồn: Lop_12_Dai_so_Chuong_6
\begin{ex}%[2D6N1-1]%[ma:59c34291]
	Nếu các biến cố $A$, $B$ thỏa mãn $P(A)>0$, $P(B)>0$ thì biểu thức $P\left(B|A\right)  $ bằng
	\choice 
		{\True $\dfrac{P\left(A|B\right)\cdot P(B)}{P(A)}$} 
		 {$\dfrac{P\left(A|B\right)\cdot P(A)}{P(B)}$}
		 {$\dfrac{P\left(A|B\right)}{P(B)}$}
		 {$\dfrac{P\left(A|B\right)}{P(A)}$}
	\loigiai{Theo lý  thuyết về xác suất.}
\end{ex}

% nguồn: Lop_12_Dai_so_Chuong_6
\begin{ex}%[2D6N1-1]%[ma:113a6ce0]
	Cho hai biến cố  $X$ và $Y$ với $\mathrm{P}(Y) > 0$. Khi đó xác suất $\mathrm{P}(X\mid Y)$ bằng
	\choice 
		 {$\dfrac{\mathrm{P}(Y)}{\mathrm{P}(X)}$}
		 {$\dfrac{\mathrm{P}(XY)}{\mathrm{P}(Y)}$}
		{\True  $\dfrac{\mathrm{P}(X)}{\mathrm{P}(Y)}$}
		 {$\dfrac{\mathrm{P}(XY)}{\mathrm{P}(X)}$}
	\loigiai{ 
		$\mathrm{P}(X\mid  Y)=\dfrac{\mathrm{P}(X)}{\mathrm{P}(Y)}$.
	 }
\end{ex}

% nguồn: Lop_12_Dai_so_Chuong_6
\begin{ex}%[2D6N1-1]%[ma:c447193b]
	Cho $A$ và $B$ là hai biến cố bất kì, với $\mathrm{P}(B)>0$.  Khi đó
	 \choice
		{$\mathrm{P}\left(B\mid A\right)=\dfrac{\mathrm{P}\left(A\cap B\right)}{\mathrm{P}(A)}$} 
		 {\True $\mathrm{P}\left(A\mid B\right)=\dfrac{\mathrm{P}\left(A\cap B\right)}{\mathrm{P}(B)}$}
		{$\mathrm{P}\left(A\mid  B\right)=\dfrac{\mathrm{P}(B)}{\mathrm{P}\left(A\cap B\right)}$}
		{$\mathrm{P}\left(B\mid A\right)=\dfrac{\mathrm{P}\left(A\cap  B\right)}{\mathrm{P}(B)}$}
	 \loigiai{
	Với  $\mathrm{P}(B)>0$, ta có $\mathrm{P}\left(A\mid B\right)=\dfrac{\mathrm{P}\left(A\cap B\right)}{\mathrm{P}(B)}$.
	} 
\end{ex}

% nguồn: Lop_12_Dai_so_Chuong_6
\begin{ex}%[2D6N1-1]%[ma:b93230bb]
	Từ một hộp có 4 tấm thẻ cùng  loại được ghi số lần lượt từ 1 đến 4. Bạn An lấy ra một cách ngẫu nhiên một thẻ từ hộp, bỏ thẻ đó ra ngoài và lại lấy một cách ngẫu nhiên thêm một thẻ nữa. Xét biến cố $A$ là \lq\lq thẻ lấy ra lần thứ nhất ghi số 3\rq\rq. Số các kết quả thuận lợi của biến cố $A$ là
	 \choice
		{\True $3$} 
		 {$2$}
		{$4$} 
		{$1$} 
	\loigiai{ 
		Tập hợp các kết quả  thuận lợi cho biến cố $A$là $\left\{\left(3;1\right),\left(3;2\right),\left(3;4\right)\right\}$.\\
	 Vậy $n\left(A\right)=3$.}
\end{ex}

% nguồn: Lop_12_Dai_so_Chuong_6
\begin{ex}%[2D6N1-1]%[ma:1f9848c3]
	Cho $A$, $B$ là $2$ biến cố bất kỳ, với $\mathrm{P}(B)>0$ và  $\mathrm{P}(AB)>0$. Khẳng định nào sau đây đúng?
	\choice 
		 {\True $\mathrm{P}(A|B)=\dfrac{\mathrm{P}(AB)}{\mathrm{P}(B)}$}
		 {$\mathrm{P}(A|B)=\dfrac{\mathrm{P}(AB)}{\mathrm{P}(A)}$}
		{$\mathrm{P}(A|B)=\dfrac{\mathrm{P}(A)}{\mathrm{P}(AB)}$} 
		 {$\mathrm{P}(A|B)=\dfrac{\mathrm{P}(B)}{\mathrm{P}(AB)}$}
	\loigiai{ 
		Theo công thức tính xác suất có điều kiện:  $\mathrm{P}(A|B)=\dfrac{\mathrm{P}(AB)}{\mathrm{P}(B)}$.
	} 
\end{ex}

% nguồn: Lop_12_Dai_so_Chuong_6
\begin{ex}%[2D6N1-1]%[ma:3945449c]
	Cho hai biến  cố $A$ và $B$ bất kì, với $\mathrm{P}(B) > 0$. Khẳng định nào sau đây là đúng?
	\choice 
		{$\mathrm{P}(A|B) = \dfrac{\mathrm{P}(A \cup B)}{\mathrm{P}(B)}$} 
		{$\mathrm{P}(A|B) = \dfrac{\mathrm{P}(A) \cdot \mathrm{P}(B)}{\mathrm{P}(B)}$} 
		{\True $\mathrm{P}(A|B)  = \dfrac{\mathrm{P}(AB)}{\mathrm{P}(B)}$}
		{$\mathrm{P}(A|B) = \dfrac{\mathrm{P}(A) +  \mathrm{P}(B)}{\mathrm{P}(B)}$}
	\loigiai{Công thức xác suất có điều  kiện là $\mathrm{P}(A|B) = \dfrac{\mathrm{P}(A \cap B)}{\mathrm{P}(B)}$.
	 }
\end{ex}

% nguồn: Lop_12_Dai_so_Chuong_6
\begin{ex}%[2D6N1-1]%[ma:b50f9142]
	Cho hai biến cố $A$ và $B$. Xác suất của biến cố $A$ với điều kiện biến cố $B$ đã xảy  ra được gọi là xác suất của $A$ với điều kiện $B$, kí hiệu là $\mathrm{P}(A|B)$. Chọn khẳng định đúng trong các khẳng định sau đây
	\choice 
		{Nếu $\mathrm{P}(A) > 0$ thì $\mathrm{P}(A|B) =  \dfrac{\mathrm{P}(A \cap B)}{\mathrm{P}(A)}$}
		{\True Nếu $\mathrm{P}(B) > 0$ thì $\mathrm{P}(A|B)  = \dfrac{\mathrm{P}(A \cap B)}{\mathrm{P}(B)}$}
		{Nếu $\mathrm{P}(B) > 0$ thì $\mathrm{P}(A|B) = \dfrac{\mathrm{P}(A \cap B)}{\mathrm{P}(A)}$} 
		{Nếu $\mathrm{P}(A) > 0$ thì $\mathrm{P}(A|B) = \dfrac{\mathrm{P}(A \cap B)}{\mathrm{P}(B)}$} 
	\loigiai{ 
		Theo định nghĩa xác suất có điều kiện: Xác suất  của biến cố $A$ với điều kiện biến cố $B$ đã xảy ra (biến cố $B$ có xác suất dương), kí hiệu $\mathrm{P}(A|B)$, được xác định bởi công thức:
		$$ \mathrm{P}(A|B) = \dfrac{\mathrm{P}(A \cap  B)}{\mathrm{P}(B)} \quad (\text{với } \mathrm{P}(B) > 0). $$
	} 
\end{ex}

% nguồn: Lop_12_Dai_so_Chuong_6
\begin{ex}%[2D6N1-1]%[ma:a06590ff]
	Một hộp chứa bốn viên bi cùng loại ghi số lần lượt từ $1$ đến $4$. Bạn Mạnh lấy ra một cách ngẫu nhiên một  viên bi, bỏ viên bi đó ra ngoài và lấy ra một cách ngẫu nhiên thêm một viên bi nữa. Không gian mẫu của phép thử đó là
	 \choice
		 {$\Omega =\left\{\left(1,2\right);\left(1,3\right);\left(1,4\right);\left(2,3\right);\left(2,4\right);\left(3,4\right)\right\}$}
		 {$\Omega =\left\{\left(1,2\right);\left(1,1\right);\left(1,3\right);\left(1,4\right);\left(2,1\right);\left(2,3\right);\left(2,4\right);\left(3,1\right);\left(3,2\right);\left(3,4\right);\left(4,1\right);\left(4,2\right);\left(4,3\right)\right\}$}
		{$\Omega  =\left\{\left(1,2\right);\left(1,3\right);\left(1,4\right);\left(2,1\right);\left(2,2\right);\left(2,3\right);\left(2,4\right);\left(1,1\right);\left(3,4\right);\left(4,4\right);\left(3,3\right)\right\}$}
		 {\True $\Omega =\left\{\left(1,2\right);\left(1,3\right);\left(1,4\right);\left(2,1\right);\left(2,3\right);\left(2,4\right);\left(3,1\right);\left(3,2\right);\left(3,4\right);\left(4,1\right);\left(4,2\right);\left(4,3\right)\right\}$}
	 \loigiai{
	 $\Omega =\left\{\left(1,2\right);\left(1,3\right);\left(1,4\right);\left(2,1\right);\left(2,3\right);\left(2,4\right);\left(3,1\right);\left(3,2\right);\left(3,4\right);\left(4,1\right);\left(4,2\right);\left(4,3\right)\right\}$.}
\end{ex}

% nguồn: Lop_12_Dai_so_Chuong_6
\begin{ex}%[2D6N1-1]%[ma:04fb8ed6]
	 Cho hai biến cố $A$ và $B$ với $\mathrm{P}(B)>0$ thì xác suất của biến cố $A$ với điều kiện biến cố $B$ đã xảy ra là
	 \choice
		 {\True $\mathrm{P}(A|B)=\dfrac{\mathrm{P}(AB)}{\mathrm{P}(B)}$}
		 {$\mathrm{P}(A|B)=\dfrac{\mathrm{P}(A)}{\mathrm{P}(B)}$}
		 {$\mathrm{P}(B|A)=\dfrac{\mathrm{P}(AB)}{\mathrm{P}(A)}$}
		{$\mathrm{P}(A|B)=\mathrm{P}(A)\mathrm{P}(B)$} 
	\loigiai{Theo công thức xác  suất có điều kiện, ta có $\mathrm{P}(A|B)=\dfrac{\mathrm{P}(AB)}{\mathrm{P}(B)}$.
	} 
\end{ex}

% nguồn: Lop_12_Dai_so_Chuong_6
\begin{ex}%[2D6N1-1]%[ma:ac3878c3]
	Cho hai biến cố $A$ và $B$ bất kì, với $\mathrm{P}(B) > 0$. Khi đó xác suất của $A$ với điều kiện $B$  là
	\choice 
		 {$\mathrm{P}(A\mid B)=\dfrac{\mathrm{P}(B)}{\mathrm{P}(A)}$}
		{$\mathrm{P}(A\mid B)=\dfrac{\mathrm{P}(A\cdot  B)}{\mathrm{P}(A)}$}
		{$\mathrm{P}(A\mid  B)=\dfrac{\mathrm{P}(A)}{\mathrm{P}(B)}$}
		 {\True $\mathrm{P}(A\mid B)=\dfrac{\mathrm{P}(A\cdot B)}{\mathrm{P}(B)}$}
	\loigiai{ 
		Với  $\mathrm{P}(B) > 0$, xác suất của $A$ với điều kiện $B$ là $\mathrm{P}(A\mid B)=\dfrac{\mathrm{P}(A\cdot B)}{\mathrm{P}(B)}$.
	} 
\end{ex}

% nguồn: Lop_12_Dai_so_Chuong_6
\begin{ex}%[2D6N1-1]%[ma:2044f99c]
	Cho $A$, $B$ là hai biến  cố của một phép thử có không gian mẫu gồm hữu hạn kết quả đồng khả năng, ký hiệu biến cố giao của $A$ và $B$ là $AB$ hay $A\cap B$, biết $\mathrm{P}(A)>0$. Xác suất của $B$ với điều kiện $A$ là
	 \choice
		 {$\mathrm{P}(B|A)=\dfrac{\mathrm{P}(A)}{\mathrm{P}(AB)}$}
		 {$\mathrm{P}(B|A)=\dfrac{\mathrm{P}(B)}{\mathrm{P}(A)}$}
		{$\mathrm{P}(B|A)=\dfrac{\mathrm{P}(AB)}{\mathrm{P}(B)}$} 
		 {\True $\mathrm{P}(B|A)=\dfrac{\mathrm{P}(AB)}{\mathrm{P}(A)}$}
	\loigiai{ 
		Công thức đúng  là $\mathrm{P}(B|A)=\dfrac{\mathrm{P}(AB)}{\mathrm{P}(A)}$.
	 }
\end{ex}

% nguồn: Lop_12_Dai_so_Chuong_6
\begin{ex}%[2D6N1-1]%[ma:7053c5bc]
	Cho $A$, $B$  là các biến cố của một phép thử $T$. Biết rằng $\mathrm{P}(B)>0$, xác suất của biến cố $A$ với điều kiện biến cố $B$ đã xảy ra được tính theo công thức nào sau đây?
	 \choice
		{$\mathrm{P}(A \mid B)=\dfrac{\mathrm{P}(AB)}{\mathrm{P}(A)  \cdot \mathrm{P}(B)}$}
		 {\True $\mathrm{P}(A \mid B)=\dfrac{\mathrm{P}(AB)}{\mathrm{P}(B)}$}
		{$\mathrm{P}(A \mid  B)=\dfrac{\mathrm{P}(A)}{\mathrm{P}(B)}$}
		{$\mathrm{P}(A  \mid B)=\dfrac{\mathrm{P}(BA)}{\mathrm{P}(A)}$}
	 \loigiai{
	Áp dụng lý thuyết ta có  $\mathrm{P}(A \mid B)=\dfrac{\mathrm{P}(AB)}{\mathrm{P}(B)}$.}
\end{ex}

% nguồn: Lop_12_Dai_so_Chuong_6
\begin{ex}%[2D6N1-1]%[ma:3f0a7f6f]
	Cho hai biến cố $A$ và $B$, trong đó $\mathrm{P}(B)>0$. Khi đó  $\mathrm{P}(A\mid B)$ bằng
	 \choice
		{$\mathrm{P}(A\cap B) \cdot  \mathrm{P}(B)$}
		{\True  $\dfrac{\mathrm{P}(A\cap B)}{\mathrm{P}(B)}$}
		{$\mathrm{P}(A\cap  B)\cdot \mathrm{P}(A)$}
		{$\dfrac{\mathrm{P}(A\cap B)}{\mathrm{P}(A)}$} 
	\loigiai{Cho hai biến cố $A$ và $B$,  trong đó $\mathrm{P}(B) > 0$. Khi đó $\mathrm{P}(A\mid B)=\dfrac{\mathrm{P}(A\cap B)}{\mathrm{P}(B)}$.}
\end{ex}

% nguồn: Lop_12_Dai_so_Chuong_6
\begin{ex}%[2D6N1-1]%[ma:8efdda38]
	Cho hai biến cố $A$ và $B$ là hai biến cố độc lập, với  $\mathrm{P}(A)=0{,}2024$; $\mathrm{P}(B)=0{,}2025$. Tính $\mathrm{P}\left(A\mid B\right)$.
	\choice 
		 {$0{,}7976$}
		{\True  $0{,}2024$}
		{$0{,}2025$} 
		{$0{,}7975$} 
	 \loigiai{
	Vì hai biến cố $A$ và $B$ là  hai biến cố độc lập nên $\mathrm{P}\left(A\mid B\right)=\mathrm{P}(A)=0{,}2024$.
	} 
\end{ex}

% nguồn: Lop_12_Dai_so_Chuong_6
\begin{ex}%[2D6N1-1]%[ma:652f714c]
	Cho $A$, $B$ là hai biến cố bất kì và $P(B)>0$. Kí  hiệu $A\cap B=AB$, công thức tính xác suất có điều kiện nào sau đây đúng?
	\choice 
		 {$P(A\mid B)=\dfrac{P(A B)}{P(A)}$}
		{$P(AB)=\dfrac{P(A\mid B)}{P(B)}$} 
		{$P(A\mid B)=P(AB)\cdot P(B)$} 
		{\True  $P(A\mid B)=\dfrac{P(AB)}{P(B)}$}
	\loigiai{Công thức đúng là $P(A\mid B)=\dfrac{P(AB)}{P(B)}$.} 
\end{ex}

% nguồn: Lop_12_Dai_so_Chuong_6
\begin{ex}%[2D6H1-1]%[ma:b4f1fdbb]
	Cho hai biến cố $A$, $B$ của cùng một phép thử thỏa $\mathrm{P}(A) = 0{,}5$; $\mathrm{P}(B) = 0{,}6$; $\mathrm{P}(A \cup B) = 0{,}7$. Khi đó  $\mathrm{P}(B|A)$ bằng
	 \choice
		 {$\dfrac{2}{3}$}
		{$0{,}6$} 
		 {$0{,}8$}
		{\True $0{,}5$} 
	\loigiai{ 
		Ta  có công thức
		\[ 
		\mathrm{P}(A \cup B) = \mathrm{P}(A) + \mathrm{P}(B)  - \mathrm{P}(A \cap B) \Rightarrow 0{,}7 = 0{,}5 + 0{,}6 - \mathrm{P}(A \cap B).
		 \]
		Suy ra $\mathrm{P}(A \cap B) = 0{,}4$. Do  đó
		\[\mathrm{P}(B|A) = \dfrac{\mathrm{P}(A \cap B)}{\mathrm{P}(A)}  = \dfrac{0{,}4}{0{,}5} = \dfrac{4}{5}=0{,}8.\]
	 }
\end{ex}


%%%%%%%%%%%%%%%% ĐÚNG/SAI %%%%%%%%%%%%%%%%

% nguồn: Lop_12_Dai_so_Chuong_6
\begin{ex}%[2D6H1-1]%[ma:caf6296e]
	Một vòng quay được chia thành $12$ phần bằng nhau và được đánh số từ $1$ đến $12$ như hình  vẽ bên dưới.
	\begin{center}
		\begin{tikzpicture}[scale=1,line join=round,line cap=round,font=\footnotesize,>=stealth]
				\def\cungelip#1#2#3#4#5{($#1+({#4*cos (#2)},{#5*sin (#2)})$) arc (#2:#3:#4  cm and #5 cm)}
				\def\diemelip#1#2#3#4{($#1+({#3*cos (#2)},{#4*sin (#2)})$)} 
				 \draw[red]\cungelip{(0,0)}{0}{360}{2}{1};
				\foreach \quay  in{0,180}{
					\begin{scope}[rotate=\quay]
					\draw[fill=red!50,](0,0)--\cungelip{(0,0)}{0}{30}{2}{1}--cycle; 
					\draw[fill=yellow!50,](0,0)--\cungelip{(0,0)}{30}{60}{2}{1}--cycle; 
					 \draw[fill=white!50,](0,0)--\cungelip{(0,0)}{60}{90}{2}{1}--cycle;
					 \draw[fill=blue!50,](0,0)--\cungelip{(0,0)}{90}{120}{2}{1}--cycle;
					\draw[fill=green!50,](0,0)--\cungelip{(0,0)}{120}{150}{2}{1}--cycle; 
					\draw[fill=black!10,](0,0)--\cungelip{(0,0)}{150}{180}{2}{1}--cycle; 
					\end{scope}
					\path 
					(0,0) coordinate  (O)
					(.7,.35) coordinate (A); 
					\draw [line width=1pt,color=red,->] (0,0)--(.7,.35); 
					\draw  (1.5,.25) node{$1$};
					\draw  (1,.5) node{$2$};
					\draw (.22,.6)  node{$3$};
					\draw (-.22,.6)  node{$4$};
					\draw (-1,.5) node{$5$}; 
					\draw (-1.5,.25) node{$6$}; 
					\draw (-1.5,-.25) node{$7$}; 
					\draw (-1,-.5)  node{$8$};
					\draw (-.22,-.6)  node{$9$};
					\draw (.22,-.6)  node{$10$};
					\draw (1,-.5) node{$11$}; 
					\draw (1.5,-.25)  node{$12$};
				 }
			\end{tikzpicture}
	\end{center}
	Xét phép  thử An và Bình lần lượt quay vòng quay trên.\\
	Gọi $A$ là biến cố \lq\lq An quay được  số chia hết cho $3$ \rq\rq; $B$ là biến cố \lq\lq An quay được số chia hết cho $5$\rq\rq; $C$ là biến cố \lq\lq Bình quay được số chẵn\rq\rq.
	 \choiceTF
		{Không gian mẫu của phép thử có số kết quả là $24$} 
		{\True Số kết quả thuận lợi cho biến cố $A$, $B$,  $C$ lần lượt là $48$, $24$, $72$}
		{Xác suất để Bình quay được số chẵn,  biết An quay được số chia hết cho $3$ là $\dfrac{1}{6}$}
		{Xác suất để An quay được số chia hết cho $5$, biết Bình quay được số lẻ là $\dfrac{1}{12}$} 
	 \loigiai{
		\begin{itemchoice}
			\itemch 
			Khi An quay thì có $12$ khả năng xảy  ra còn khi Bình quay thì có $12$ khả năng xảy ra.\\
			Theo quy  tắc nhân có $12\cdot 12=144$ khả năng xảy ra nên $n\left(\Omega\right) =144$.
			 \itemch
			Xét biến cố $A$, An quay được số chia hết cho $3$ thì có $4$ khả năng còn khi Bình quay một số ngẫu  nhiên có $12$ khả năng.\\
			Theo quy tắc nhân có  $4\cdot 12=48$ khả năng nên $n\left(A\right)=48$.\\
			Xét biến cố $B$, An quay được số chia hết cho $5$ có $2$ khả năng.\\ 
			Bình quay một số ngẫu nhiên có $12$ khả năng nên theo quy tắc nhân có $2\cdot 12=24$ khả năng. Do  đó $n\left(B\right)=24$.\\
			Xét biến cố $C$, An quay một số ngẫu nhiên có $12$ khả năng. Bình quay một số chẵn có $6$ khả năng nên theo quy tắc nhân có $12\cdot 6=72$  khả năng. Do đó $n\left(C\right)=72$.
			\itemch 
			Khi An quay được số  chia hết cho $3$ thì số khả năng Bình quay được số chẵn là $6$ khả năng.\\
			Do đó xác suất để Bình quay được số chẵn  biết An quay được số chia hết cho $3$ là $\dfrac{6}{12}=\dfrac{1}{2}$.
			\itemch 
			Khi Bình quay được số lẻ thì số khả năng An quay được số chia hết cho $5$  là $2$ khả năng. \\
			Do đó xác suất để An quay được số chia hết  cho $5$, biết Bình quay được số lẻ là $\dfrac{2}{12}=\dfrac{1}{6}$.
		\end{itemchoice}
	} 
\end{ex}


%%%%%%%%%%%%%%%% TRẢ LỜI NGẮN %%%%%%%%%%%%%%%%

% nguồn: Lop_12_Dai_so_Chuong_6
\begin{bt}%[2D6H1-1]%[ma:34e39e84]
	Một nhóm học sinh có $5$ nam và $7$ nữ tham gia lao động trên sân trường. Cô giáo chọn ngẫu nhiên đồng thời $3$ bạn trong nhóm đi quét sân. Biết rằng có ít nhất một bạn nam được chọn, tính xác suất để ba bạn được chọn có cùng giới tính \textit{(kết quả làm  tròn đến hàng phần trăm)}.
	\shortans{0,05} 
	 \loigiai{
		Số cách chọn $3$ học sinh trong đó có ít nhất một bạn nam là  $\mathrm{C}_{12}^{3}-\mathrm{C}_{7}^{3}=185$. \\
		Số cách chọn $3$ học sinh cùng giới tính nam là  $\mathrm{C}_{5}^{3}=10$. \\
		Vậy xác suất để ba bạn được chọn có cùng giới  tính biết có ít nhất một bạn nam được chọn là $\dfrac{10}{185} \approx 0{,}05$.
	 }
\end{bt}


%%%%%%%%%%%%%%%% TỰ LUẬN %%%%%%%%%%%%%%%%

% nguồn: Lop_12_Dai_so_Chuong_6
\begin{bt}%[2D6H1-1]%[ma:45ee08f8]
	Chứng tỏ rằng nếu $A$ và $B$ là hai biến cố độc lập thì $\mathrm{P}(\overline{A}\mid B) = \mathrm{P}(\overline{A})$ và $\mathrm{P}(A\mid \overline{B}) = \mathrm{P}(A)$. 
	\loigiai{ 
		Ta có $\mathrm{P}(\overline{A} \mid B)$ là xác suất của $\overline{A}$ biết rằng biến cố $B$ đã xảy ra. Vì $A$ và $B$ là hai biến cố độc lập nên $\overline{A}$  và $B$ độc lập, tức là việc xảy ra $B$ không ảnh hưởng đến xác suất không xuất hiện của $A$. Do đó $\mathrm{P}(\overline{A} \mid B) = \mathrm{P}(\overline{A})$.\\
		Chứng minh $\mathrm{P}(A \mid \overline{B})  = \mathrm{P}(A)$.\\
		Tương tự $\mathrm{P}(A \mid \overline{B})$ là xác suất của $A$ biết rằng biến cố $\overline{B}$ đã xảy ra. Vì $A$ và $B$ là hai biến cố độc lập nên $A$ và $\overline{B}$ độc  lập, tức là việc không xảy ra $B$ không ảnh hưởng đến xác suất xuất hiện của $A$. Do đó $\mathrm{P}(A \mid \overline{B}) = \mathrm{P}(A)$.
	 }
\end{bt}

% nguồn: Lop_12_Dai_so_Chuong_6
\begin{bt}%[2D6H1-1]%[ma:965fa372]
	Có hai hộp chứa các  thẻ được đánh số. Hộp thứ nhất có các thẻ được đánh số từ 1 đến 4, hộp thứ hai có các thẻ được đánh số từ 5 đến 6. Các thẻ có cùng kích thước và khối lượng. Bạn Phương lấy ngẫu nhiên một thẻ từ hộp thứ nhất bỏ vào hộp thứ hai. Sau đó bạn lại lấy ngẫu nhiên một thẻ từ hộp thứ hai. Liệt kê các kết quả của phép thử biết lần thứ nhất bạn Phương lấy được một thẻ đánh số chẵn.
	\loigiai{ 
		Vì đã biết lần  thứ nhất bạn Phương lấy được một thẻ đánh số chẵn, nghĩa là có thể lấy được thẻ đánh số $2$ hoặc $4$.
		\begin{itemize}
			\item Nếu bạn Phương lấy được thẻ đánh số $2$ và bỏ vào  hộp thứ hai, thì lúc này trong hộp thứ hai có các thẻ đánh số $2, 5, 6$. Khi đó các kết quả có thể là: $(2;2), (2;5), (2;6)$.
			\item Nếu bạn Phương lấy được thẻ đánh số $4$ và bỏ vào hộp thứ hai, thì lúc này trong hộp thứ hai có các thẻ đánh số $4, 5,  6$. Khi đó các kết quả có thể là: $(4;4), (4;5), (4;6)$.
		\end{itemize}
		Vậy các kết quả của phép thử là: $(2;2), (2;5), (2;6), (4;4),  (4;5), (4;6)$.
	} 
\end{bt}
